\documentclass[t]{beamer}

% For more themes, color themes and font themes, see:
% http://deic.uab.es/~iblanes/beamer_gallery/index_by_theme.html
%\usetheme{iclpt}
\usepackage{ragged2e}
\usepackage{etoolbox}
\apptocmd{\frame}{}{\justifying}{}

\apptocmd{\column}{}{\justifying}{}
\hyphenpenalty=1000
\exhyphenpenalty=1000

%
\mode<presentation>
{
  \usetheme{Madrid}       % Madrid or try default, Darmstadt, Warsaw, ...
  \usecolortheme{spruce} % default or try albatross, beaver, crane, beetle, dove, fly, monarca, seagull, spruce, and wolverine ...
  \usefonttheme{serif}    % serif or try default, structurebold, ...
  %\setbeamertemplate{navigation symbols}{}
  \setbeamertemplate{caption}[numbered]
} 

\usepackage[english]{babel}
\usepackage[utf8x]{inputenc}
\usepackage{chemfig}
\usepackage[version=3]{mhchem}

\usepackage{graphics}
\usepackage{graphicx}
\usepackage{pifont}
\usepackage{amsfonts}
\usepackage{amssymb}
\usepackage{xspace,colortbl}
%\usepackage[screen,panelleft,code,blue,sectionbreak]{pdfscreen}
\usepackage{utopia}
\usepackage[round]{natbib}
\usepackage{hyperref}
\usepackage{lmodern}
\usepackage[T1]{fontenc}
% On Overleaf, these lines give you sharper preview images.
% You might want to `comment them out before you export, though.
\usepackage{pgfpages}
\pgfpagesuselayout{resize to}[%
  physical paper width=8in, physical paper height=6in]
  
\newtheorem{thm}{Theorem}[section]
\newtheorem{cor}[thm]{Corollary}
\newtheorem{lem}[thm]{Lemma}
\newtheorem{prop}[thm]{Proposition}
%\theoremstyle{definition}
\newtheorem{defn}[thm]{Definition}
%\theoremstyle{remark}
\newtheorem{rem}[thm]{Remark}
\numberwithin{equation}{section}
% MATH -----------------------------------------------------------
\newcommand{\norm}[1]{\left\Vert#1\right\Vert}
\newcommand{\abs}[1]{\left\vert#1\right\vert}
\newcommand{\set}[1]{\left\{#1\right\}}
\newcommand{\Real}{\mathbb R}
\newcommand{\eps}{\varepsilon}
\newcommand{\To}{\longrightarrow}
\newcommand{\BX}{\mathbf{B}(X)}
\newcommand{\A}{\mathcal{A}}
%\usetheme{iclpt}
\usepackage{url}
\usepackage{amsmath}


\babelprovide[import,
  onchar = fonts ids,
  hyphenrules = +, % Use empty patters
  ]{amharic}

\babelposthyphenation{amharic}{([ሀሁሂሃሄህሆለሉሊላሌልሎሏሐሑሒሓሔሕሖሗመሙሚማሜምሞሟሠሡሢሣሤሥሦሧረሩሪራሬርሮሯሰሱሲሳሴስሶሷሸሹሺሻሼሽሾሿቀቁቂቃቄቅቆቈቊቋቌቍበቡቢባቤብቦቧቨቩቪቫቬቭቮቯተቱቲታቴትቶቷቸቹቺቻቼችቾቿኀኁኂኃኄኅኆኈኊኋኌኍነኑኒናኔንኖኗኘኙኚኛኜኝኞኟአኡኢኣኤእኦኧከኩኪካኬክኮኰኲኳኴኵኸኹኺኻኼኽኾወዉዊዋዌውዎዐዑዒዓዔዕዖዘዙዚዛዜዝዞዟዠዡዢዣዤዥዦዧየዩዪያዬይዮደዱዲዳዴጼጽጾጿፀፁፂፃፄፅፆፈፉፊፋፌፍፎፏፐፑፒፓፔፕፖፗ])}{
 { no = {1}, post = {1} }
}

\babelfont[amharic]{rm}{FreeSerif}

\usepackage{newunicodechar}

\newunicodechar{፡}{፡\hskip0pt plus 3pt\relax}
\newunicodechar{።}{።\hskip0pt plus 3pt\relax}
%\pgfdeclareimage[height=0.5cm]{university-logo}{eiwr.png}
% \logo{\pgfuseimage{eiwr.png}}
% Here's where the presentation starts, with the info for the title slide
\begin{block}{}
\begin{flushleft}
\begin{figure}
\includegraphics[scale=0.35]{Belha}
\end{figure}
\end{flushleft}
\end{block}
\title[በማሞ ካሠኝ]{ተግባራዊ ስልጠና: በውኃና አየር ንብረት መረጃ ላይ የተመሠረተ የኮምፒውተር ስልጠና መርሐ ግብር \\ (R, Python, MATLAB)}

\author[በልሃ ልበልሻ የምርምር ተቋም]{በማሞ ካሠኝ \\ \url mamo.kassegn@aau.edu.et}



\date{\today}
\vfill
\begin{document}
\titlepage 

% These three lines create an automatically generated table of contents.
\newpage
\begin{frame}{ይዘት}
\vspace{-6mm}
\small{
 ~~~~~~~~~~~ \tableofcontents
 } \tiny
\section{የኮምፒውተር መተየቢያ ፕሮግራሚንግ }
\subsection{የኮምፒውተር መትየቢያ አጫጫን}
\subsection{መሠረታዊ የኮምፒውተር ፕሮግራሚንግ ዕውቀት}
\subsection{የተሰበሰቡ መረጃዎችን በኮምፒውተር የመክፈት ስልት(CSV, text, NetCDF) }
\subsection{የስታቲስቲካልና የስዕላዊ የመረጃ አገላለፅ ዘዴዎች}
\subsection{የጎደሉ መረጃዎችን በአርቲፊሻል ጥምርታዎች የመሙላት ስልት}
\subsection{ዓመታዊና ወቅታዊ የዝናብና የወንዞች ፍሰት መስመራዊ መግለጫ}
\subsection{የየጊዜውን ሁኔታ የሚገልጥ የዝናብና የወንዞች ፍሰት መስመር (PDC \& FDC)}
\section{ብርቱ የኮምፒውተር ኮዲንግ ትግበራ}
\subsection{በረዥም ጊዜ የተሰበሰቡ የዝናብ፤ የሙቀትና ሌሎች የአየር ንብረት መረጃዎች}
\subsection{ጅኦስታቲስቲካል ስንሰላ }
\section{netCDF ፋይሎችን ማንበብ}
\subsection{ዘንዶው ስለተባለው የኮምፒውተር ፕሮግራሚንግ መተግበሪያ}
\subsection{ዘንዶ የኮምፒውተር ማሰናሰያና የአመሠራረት ዘይቤ} 
\subsection{ገቢና ወጪ ፋይሎች}
\subsection{በኮምፒውተር ተንቀሳቃሽ ምስሎችን }
\end{frame}
\begin{frame}{የኮምፒውተር ፕሮግራሚንግ መሠረታዊ }
\justifying{
\begin{itemize}


\item የኮምፒዩተር ፕሮግራሚንግ የተወሰኑ ስሌቶችን (ወይም በአጠቃላይ ፣ የተወሰኑ የኮምፒዩተር ውጤቶችን በማከናወን) ብዙውን ጊዜ ተፈጻሚ የሚሆኑ የኮምፒተር ፕሮግራሞችን የመንደፍ እና የመገንባት ሂደት ነው። 

\item ፕሮግራሚንግ እንደ ትንተና፣ አልጎሪዝም ማመንጨት፣ የመገለጫ ስልተ ቀመሮችን ትክክለኛነት እና የሃብት ፍጆታ እና የአልጎሪዝም አተገባበርን (በተለምዶ በተለየ የፕሮግራም አወጣጥ ቋንቋ፣ በተለምዶ ኮድ ማድረግ ተብሎ የሚጠራ) ተግባራትን ያካትታል። 

\item ፕሮግራም ማዘጋጀት አንድ ኮምፒተር እንዴት እንደሚሰራ የሚያስተምር የፈጠራ ሥራ ነው። ኮምፒተር ውስጥ ቁጭ ብሎ በሴኮንዶች ውስጥ ማናቸውንም ቁልፎች በመሰደርና ቃሎችን በማምረት የድሪጊቶችን ገፅታ መቀየር ያስችላል።
\end{itemize}}
\end{frame}
\newpage
\begin{frame}
\vspace{3cm}
\centering
\Huge
አመሰግናለሁ!!
\end{frame}
\end{document}